\section{One support, both phases}\label{sec:ppt3-coexistence}

The product curves constrain biseparable decompositions, but the support alone does not decide GME. The support
$R=R_1$ of the first family carries both biseparable and GME fully-PPT states with the same rank profile.
An explicit line crosses their boundary while remaining strictly PPT. The proofs, with the exact matrices and minors, are in
Appendix~\ref{sec:ppt3-app-coexistence}.

Biseparable states on $R$ can be written down explicitly from the product curve of Proposition~\ref{prop:ppt3-witness}. For
finite $z$ let $p_z=Bq_1(z)/\|Bq_1(z)\|$, a normalised $A|BC$-product vector in $R$, let $\Pi$ be the cyclic qubit permutation,
and define the cyclically and parity symmetrised mixture
\begin{equation}\label{eq:ppt3-gamma}
 \Gamma_z=\frac16\sum_{j=0}^2\sum_{\epsilon=\pm1}
   \Pi^j\ket{p_{\epsilon z}}\bra{p_{\epsilon z}}
   \Pi^{-j}.
\end{equation}
Since $Z^{\otimes3}\ket{p_z}=-\ket{p_{-z}}$, the sum over $\epsilon=\pm1$ is the global parity average, and each $\Gamma_z$
is an equal mixture of six pure product states across the three cuts, so $\Gamma_z$ and its convex mixtures are biseparable and supported on $R$.

\begin{theorem}[biseparable and GME states on one PPT support]
\label{thm:ppt3-phase-coexistence}
\cls{proved; exact rational arithmetic, PG-2}
The fixed support $R=R_1$ contains both a
biseparable state and a GME state of rank five,
\emph{each} PPT with strictly positive partial
transposes. They therefore have the same rank
profile $(5,8,8,8)$ and exactly the same
three cut-product curves. An explicit
biseparable example is
$\sigma_{\rm BS}=\frac13\bigl(\ket{111}\bra{111}+\Gamma_2+\Gamma_{2i}\bigr)=BX_{\rm BS}B^\dagger$,
where the rational coefficient matrix $X_{\rm BS}$ in the basis $(u,v,w,e,o)$ is given in
\eqref{eq:ppt3-bs-same-support} of Appendix~\ref{sec:ppt3-app-coexistence}.
The GME example is $\rho_{10/3}$ from
Theorem~\ref{thm:ppt3-family}.
Within the trace-one operators supported on $R$,
the fully-PPT biseparable states and the fully-PPT
GME states both have nonempty relative interiors.
\end{theorem}

\emph{Idea of the proof.} The state $\sigma_{\rm BS}$ is biseparable by construction; expanding the mixture gives
$X_{\rm BS}$, and Sylvester's criterion on the two parity blocks of one partial transpose gives strict PPT. The state
$\rho_{10/3}$ has the same range and is GME by \eqref{eq:ppt3-h5-value}. Strict PPT, rank five and GME persist under small
perturbations within $R$. On the biseparable side, averaging curve projectors over a measure of full support gives an
interior point of the supported biseparable cone: no nonzero supporting form can vanish on that average. Mixing it into
$\sigma_{\rm BS}$ retains strict PPT (Appendix~\ref{sec:ppt3-app-coexistence}).

\begin{proposition}[an exposed phase boundary inside the strict-PPT region]
\label{prop:ppt3-exposed-transition}
\cls{proved; exact rational arithmetic, PG-4}
With the biseparable cut-product mixtures $\Gamma_z$ of
\eqref{eq:ppt3-gamma}, define
\begin{equation}\label{eq:ppt3-exposed-state}
 \sigma_\partial=\frac15 \Gamma_1+\frac12 \Gamma_i+
                  \frac3{10}\Gamma_{i\sqrt5}
 =BX_\partial B^\dagger ,
\end{equation}
with the rational matrix $X_\partial$ of \eqref{eq:ppt3-xdelta} in Appendix~\ref{sec:ppt3-app-coexistence}.
This rank-five biseparable state has
$\sigma_\partial^{T_c}\succ I/300$ at each cut, yet
$\tr(H_5X_\partial)=0$. The straight line
\begin{equation}\label{eq:ppt3-exact-phase-line}
 \sigma(t)=(1-t)\sigma_\partial+t\sigma_{\rm BS},
          \qquad -\frac1{200}\le t\le1,
\end{equation}
consists entirely of rank-five fully-PPT states
with strictly positive partial transposes. Along
this \emph{single fixed-support, fixed-rank} line,
\[
 \sigma(t)\text{ is biseparable for }0\le t\le1,
 \qquad
 \sigma(t)\text{ is GME for }-1/200\le t<0.
\]
Thus $t=0$ is an exact biseparable--GME transition
\emph{inside} the strict-PPT region, exposed by $H_5$.
\end{proposition}

\emph{Idea of the proof.} The factorisation \eqref{eq:ppt3-h5} vanishes at $|z|=1$ and at $z=\pm i\sqrt5$, so $H_5$ has zero
expectation on each summand of $\sigma_\partial$. Since $\tr(H_5X_{\rm BS})>0$, the witness value is linear in $t$ along
the line and negative exactly for $t<0$, which excludes biseparability there. The PPT margin is a Sylvester check, a
Hilbert--Schmidt estimate keeps it for the small negative range of $t$, and convexity covers $0\le t\le1$. The exact data are
in Appendix~\ref{sec:ppt3-app-coexistence}.
